\documentclass[11pt,a4paper]{article}
\usepackage[utf8]{inputenc}
\usepackage[T1]{fontenc}
\usepackage[english]{babel}
\usepackage{amsmath,amssymb,amsthm}
\\usepackage[a4paper,margin=0.65in]{geometry}
\usepackage{pdfpages}
\emergencystretch=3em
\tolerance=600
\usepackage{graphicx}
\usepackage{array}
\usepackage{longtable}
\usepackage{booktabs}
\usepackage{fancyhdr}
\usepackage[protrusion=true,expansion=false]{microtype}
\usepackage{hyperref}

% Fallbacks for non-standard commands
\providecommand{\half}{\tfrac{1}{2}}
\providecommand{\third}{\tfrac{1}{3}}
\providecommand{\GF}{\mathrm{GF}}
% Cayley uses an inverted Greek omega for the totient (Euler's phi):
\providecommand{\totient}{\varpi}
% Cayley's "sum of divisors" symbol -- here we use the integral sign
% (as it appears in the Cambridge text), via a dedicated macro:
\providecommand{\sumdiv}{{\textstyle\int}}

\pagestyle{fancy}
\fancyhf{}
\rhead{Cayley --- Collected Papers, Vol.~IX}
\lhead{Pages 456--478}
\cfoot{\thepage}

\title{\textbf{Arthur Cayley}\\\vspace{0.5em}
\Large Collected Mathematical Papers\\\vspace{0.3em}
\large Volume IX, Pages 456--478 (Papers 610 conclusion \& 611)}
\author{Transcribed from the Original Publications}
\date{}

\begin{document}
\maketitle
\thispagestyle{empty}

\tableofcontents
\newpage

%% ========================================================================
%% PAPER [610] -- continued and concluded, pp. 456--460 (PDF pp. 476--480)
%% ========================================================================
\section*{\textbf{610.} On the Analytical Forms called Trees, with Application to the Theory of Chemical Combinations \textit{(concluded)}}
\addcontentsline{toc}{section}{\textbf{610.} On the Analytical Forms called Trees \textit{(concluded, pp.~456--460)}}

\textit{[Conclusion of the paper begun at p.~427.   Pages 456--458 give Table VIII (Carbon Centre- and Bicentre-Trees) and the Subsidiary Table for the generating functions of Tables VII and VIII.   Pages 459--460 give Tables A and B (general valency tables) and the explanatory text leading to the chemical interpretation.]}

\bigskip

\noindent\textbf{[p.~456]}\quad

\begin{center}
\textbf{Table VIII.\ ---Carbon Centre- and Bicentre-Trees.}
\end{center}

\noindent\textit{Column headings: ``Index of $t$, or number of main branches'' and ``Index of $x$, or number of knots''.   For each value of the number of knots, the entries in the upper portion (Centre-Trees) are arranged by altitude (columns 1\,\dots\,6) and the row marked ``Total'' shows the grand total for that knot count.   The right-hand portion (Bicentre-Trees) is arranged by altitude (columns 1\,\dots\,5) and ``1'' (the bicentre column).}

\medskip

{\small
\renewcommand{\arraystretch}{1.1}
\begin{longtable}{c c | c c c c c c | c c c | c c c c c}
\hline
\multicolumn{2}{c|}{} & \multicolumn{6}{c|}{Centre-Trees (altitude)} & \multicolumn{3}{c|}{} & \multicolumn{5}{c}{Bicentre-Trees (altitude)} \\
$n$ & $a$ & 1 & 2 & 3 & 4 & 5 & 6 & Cent. & G.T. & Bic. & 1 & 2 & 3 & 4 & 5 \\
\hline
\endfirsthead
\hline
$n$ & $a$ & 1 & 2 & 3 & 4 & 5 & 6 & Cent. & G.T. & Bic. & 1 & 2 & 3 & 4 & 5 \\
\hline
\endhead
1 & 0 & 1 & & & & & & 1 & 1 & 0 & & & & & \\
\textit{Total} & & 1 & & & & & & 1 & 1 & 0 & & & & & \\
2 &   &   & & & & & & 0 & 1 & 1 & 1 & & & & \\
3 & 2 & 1 & 1 & & & & & 1 & 1 & 0 & & & & & \\
\textit{Total} & & 1 & 1 & & & & & 1 & 1 & 0 & & & & & \\
4 & 2 & 1 & 1 & & & & & 1 & 2 & 1 & 1 & & & & \\
\textit{Total} & & 1 & 1 & & & & & 1 & 2 & 1 & 1 & & & & \\
5 & 2 & 1 & & 2 & & & & 2 & 3 & 1 & 1 & & & & \\
\textit{Total} & & 1 & & 2 & & & & 2 & 3 & 1 & 1 & & & & \\
6 & 2 & 1 & 1 & & & & & & & & & & & & \\
6 & 3 & 1 & 1 & & & & & & & & & & & & \\
\textit{Total} & & 2 & 2 & & & & & 2 & 5 & 3 & 2 & 1 & & & \\
7 & 2 & & 2 & 1 & 3 & & & & & & & & & & \\
7 & 3 & & 2 &   & 2 & & & & & & & & & & \\
7 & 4 & & 1 &   & 1 & & & & & & & & & & \\
\textit{Total} & & & 5 & 1 & 6 & & & 6 & 9 & 3 & 1 & 2 & & & \\
8 & 2 & & 1 & 2 & 3 & & & & & & & & & & \\
8 & 3 & & 3 & 1 & 4 & & & & & & & & & & \\
8 & 4 & & 2 &   & 2 & & & & & & & & & & \\
\textit{Total} & & & 6 & 3 & 9 & & & 9 & 18 & 9 & 1 & 7 & 1 & & \\
9 & 2 & & 1 & 7 & 1 & & & & & & & & & & \\
9 & 3 & & 3 & 3 &   & & & & & & & & & & \\
9 & 4 & & 4 & 1 &   & & & & & & & & & & \\
\textit{Total} & & & 8 & 11 & 1 & 20 & & 20 & 35 & 15 & & 12 & 3 & & \\
10 & 2 & & & 12 & 3 & 15 & & & & & & & & & \\
10 & 3 & & 3 & 11 & 1 & 15 & & & & & & & & & \\
10 & 4 & & 4 & 3 &   & 7  & & & & & & & & & \\
\textit{Total} & & & 7 & 26 & 4 & 37 & & 37 & 75 & 38 & & 23 & 14 & 1 & \\
11 & 2 & & & 23 & 14 & 1 & 38 & & & & & & & & \\
11 & 3 & & 2 & 24 & 4 &   & 30 & & & & & & & & \\
11 & 4 & & 5 & 12 & 1 &   & 18 & & & & & & & & \\
\textit{Total} & & & 7 & 59 & 19 & 1 & 86 & 86 & 159 & 73 & & 30 & 39 & 4 & \\
12 & 2 & & & 30 & 39 & 5 & 74 & & & & & & & & \\
12 & 3 & & 1 & 54 & 19 & 1 & 75 & & & & & & & & \\
12 & 4 & & 4 & 27 & 4 &   & 35 & & & & & & & & \\
\textit{Total} & & & 5 & 111 & 62 & 5 & 183 & 183 & 357 & 174 & & 42 & 108 & 23 & 1 \\
13 & 2 & & & 42 & 108 & 23 & 1 & 174 & & & & & & & \\
13 & 3 & & 1 & 88 & 63 & 5 &   & 157 & & & & & & & \\
13 & 4 & & 4 & 63 & 20 & 1 &   &  88 & & & & & & & \\
\textit{Total} & & & 5 & 193 & 191 & 29 & 1 & 419 & 799 & 380 & & 47 & 244 & 84 & 5 \\
\hline
\end{longtable}
}

\bigskip
\noindent\textbf{[p.~457]}\quad

\begin{center}
\textbf{Subsidiary Table for $\GF$'s of Tables VII.\ and VIII.}
\end{center}

\noindent\textit{The table is arranged by ``column'' (i.e.\ the number which is in successive use for the analytical work of Tables VII and VIII).   For each column there is a generating function (``GF''), and that GF is built up from a ``first factor'' and a ``second factor''.   Within each first and second factor the rows are indexed by the exponent of $t$ ($=1,2,3,\ldots$, the number of main branches), and the columns are indexed by the exponent of $x$ ($=1,2,\ldots,13$, the number of knots).   The entries are the coefficients.}

\smallskip
\noindent The table extends across pp.~457--458 (\textit{continued}).   We reproduce the salient lines, preserving Cayley's classification of GF, first factor, second factor for each ``column'' value.

\medskip

\noindent\textit{GF, column 0:}\quad $1$ (a single entry under index $x=1$).

\noindent\textit{GF, column 1:}\quad $-1$ (under index $x=1$).

\noindent\textit{First factor (column 1, label $(1)$):}\ \ row $t=1$: $1$ at $x=1$.

\noindent\textit{Second factor (column 1):}\ \ row $t=1$: $1$ at $x=2$; $1$ at $x=3$; $1$ at $x=4$.

\medskip

\noindent\textit{GF, column 2:}\ \ $-1,-1,-1$ under indices $x=3,4,5$.

\noindent\textit{First factor (column 2):}\ \ row $t=1$: $1,1,1$ at $x=3,4,5$;\ \ row $t=2$: $1,1,2,1,1$ at $x=5,6,7,8,9$;\ \ row $t=3$: $1,1,2,2,2,1,1$ at $x=7,8,9,10,11,12,13$;\ \ row $t=4$: $1,1,2,2,3$ at $x=9,10,11,12,13$.

\noindent\textit{Second factor (column 2):}\ \ row $t=1$: $1$ at $x=2$;\ \ row $t=2$: $1$ at $x=3$;\ \ row $t=3$: $1$ at $x=4$.

\medskip

\noindent\textit{GF, column 3:}\ \ $-1,-2,-4,-4,-6,-4,-4,-3,-2,-1,-1$ under indices $x=3,4,5,6,7,8,9,10,11,12,13$.

\noindent\textit{First factor (column 3):}\ \ row $t=1$: $1,2,4,4,5,4,4,3,2,1$ at $x=3,\ldots,12$;\ \ row $t=2$: $1,2,4,12,23,30,42$ at $x=7,\ldots,13$;\ \ row $t=3$: $1,2,7,16$ at $x=10,\ldots,13$;\ \ row $t=4$: $1$ at $x=13$.

\noindent\textit{Second factor (column 3):}\ \ row $t=1$: $1,1,1,1$ at $x=2,3,4,5$;\ \ row $t=2$: $1,1,2,2,2,1,1$ at $x=3,4,5,6,7,8,9$;\ \ row $t=3$: $1,1,2,3,3,3,3,2,1,1$ at $x=4,\ldots,13$.

\medskip

\noindent\textit{GF, column 4:}\ \ $-1,-3,-8,-15,-27,-43,-67,-97,-136,-183$ under indices $x=4,5,\ldots,13$.

\noindent\textit{First factor (column 4):}\ \ row $t=1$: $1,3,8,15,27,43,67,97,136$ at $x=4,\ldots,12$;\ \ row $t=2$: $1,3,14,39,108$ at $x=9,\ldots,13$;\ \ row $t=3$: $1$ at $x=13$.

\noindent\textit{Second factor (column 4):}\ \ row $t=1$: $1,1,2,3,4,4,5,4,4,3,2,1$ at $x=2,\ldots,13$;\ \ row $t=2$: $1,1,3,5,10,14,23,29,40,46,55$ at $x=3,\ldots,13$;\ \ row $t=3$: $1,3,6,12,20,37,56,89,128$ at $x=4,5,\ldots,13$ (with a gap at $x=5$).

\medskip

\noindent\textit{GF, column 5:}\ \ $-1,-4,-13,-32,-74,-155,-316,-612,-1160$ under indices $x=5,\ldots,13$.

\noindent\textit{First factor (column 5):}\ \ row $t=1$: $4,13,32,74,155,316,612$ at $x=6,\ldots,12$;\ \ row $t=2$: $1,4,23$ at $x=11,12,13$.

\noindent\textit{Second factor (column 5):}\ \ row $t=1$: $1,1,2,7,12,20,31,47,70,99,137$ at $x=2,\ldots,13$ (with a gap);\ \ row $t=2$: $1,1,3,6,14,27,56,103,194,343,605$ at $x=3,\ldots,13$ (with a gap);\ \ row $t=3$: $1,1,3,7,16,34,76,151,307,602$ at $x=4,\ldots,13$.

\medskip

\noindent\textit{GF, column 6:}\ \ $-1,-6,-19,-56,-151,-374,-889,-2032$ under indices $x=6,\ldots,13$.

\noindent\textit{First factor (column 6):}\ \ row $t=1$: $1,5,19,56,151,374,889$ at $x=7,\ldots,13$;\ \ row $t=2$: $1$ at $x=13$.

\noindent\textit{Second factor (column 6):}\ \ row $t=1$: $1,1,2,4,8,16,33,63,121,225,415,749$ at $x=2,\ldots,13$;\ \ row $t=2$: $1,1,3,6,15,32,75,160,350,732,1534$ at $x=3,\ldots,13$;\ \ row $t=3$: $1,1,3,7,17,39,95,214,491,1093$ at $x=4,\ldots,13$.

\bigskip
\noindent\textbf{[p.~458]}\quad\textit{Subsidiary Table for $\GF$'s of Tables VII.\ and VIII.\ (continued).}

\medskip

\noindent\textit{GF, column 7:}\ \ $-1,-6,-26,-88,-267,-743,-1968$ under indices $x=7,\ldots,13$.

\noindent\textit{First factor (column 7):}\ \ row $t=1$: $1,6,26,88,267,743$ at $x=8,\ldots,13$.

\noindent\textit{Second factor (column 7):}\ \ row $t=1$: $1,1,2,4,8,17,38,82,177,376,789,1638$ at $x=2,\ldots,13$;\ \ row $t=2$: $1,1,3,6,15,33,81,186,439,1005,2304$ at $x=3,\ldots,13$;\ \ row $t=3$: $1,1,3,7,17,40,101,241,587,1402$ at $x=4,\ldots,13$.

\medskip

\noindent\textit{GF, column 8:}\ \ $-1,-7,-34,-129,-432,-1320$ under indices $x=8,\ldots,13$.

\noindent\textit{First factor (column 8):}\ \ row $t=1$: $1,7,34,129,432$ at $x=9,\ldots,13$.

\noindent\textit{Second factor (column 8):}\ \ row $t=1$: $1,1,2,4,8,17,39,88,203,464,1056,2381$ at $x=2,\ldots,13$;\ \ row $t=2$: $1,1,3,6,15,33,82,193,473,(\ldots),2743$ at $x=3,\ldots,13$;\ \ row $t=3$: $1,1,3,7,17,40,102,248,622,1540$ at $x=4,\ldots,13$.

\medskip

\noindent\textit{GF, column 9:}\ \ $-1,-8,-43,-180,-657$ under indices $x=9,\ldots,13$.

\noindent\textit{First factor (column 9):}\ \ row $t=1$: $1,8,43,180$ at $x=10,\ldots,13$.

\noindent\textit{Second factor (column 9):}\ \ row $t=1$: $1,1,2,4,8,17,39,89,210,498,1185,2813$ at $x=2,\ldots,13$;\ \ row $t=2$: $1,1,3,6,15,33,82,194,481,1178,2924$ at $x=3,\ldots,13$;\ \ row $t=3$: $1,1,3,7,17,40,102,249,630,1584$ at $x=4,\ldots,13$.

\medskip

\noindent\textit{GF, column 10:}\ \ $-1,-9,-53,-242$ under indices $x=10,\ldots,13$.

\noindent\textit{First factor (column 10):}\ \ row $t=1$: $1,9,53$ at $x=11,12,13$.

\noindent\textit{Second factor (column 10):}\ \ row $t=1$: $1,1,2,4,8,17,39,89,211,506,1228,2993$ at $x=2,\ldots,13$;\ \ row $t=2$: $1,1,3,6,15,33,82,194,482,1187,2977$ at $x=3,\ldots,13$;\ \ row $t=3$: $1,1,3,7,17,40,102,249,631,1593$ at $x=4,\ldots,13$.

\medskip

\noindent\textit{GF, column 11:}\ \ $-1,-10,-63$ under indices $x=11,12,13$.

\noindent\textit{First factor (column 11):}\ \ row $t=1$: $1,10,63$ at $x=12,13$ (and a leading entry).

\noindent\textit{Second factor (column 11):}\ \ row $t=1$: $1,1,2,4,8,17,39,89,211,507,1237,3048$ at $x=2,\ldots,13$;\ \ row $t=2$: $1,1,3,6,15,33,82,194,482,1188,2987$ at $x=3,\ldots,13$;\ \ row $t=3$: $1,1,3,7,17,40,102,249,631,1594$ at $x=4,\ldots,13$.

\medskip

\noindent\textit{GF, column 12:}\ \ $-1,-11$ under indices $x=12,13$.

\noindent\textit{First factor (column 12):}\ \ row $t=1$: $1,11$ at $x=12,13$.

\noindent\textit{Second factor (column 12):}\ \ row $t=1$: $1,1,2,4,8,17,39,89,211,507,1238,3055$ at $x=2,\ldots,13$;\ \ row $t=2$: $1,1,3,6,15,33,82,194,482,1188,2988$ at $x=3,\ldots,13$;\ \ row $t=3$: $1,1,3,7,17,40,102,249,631,1594$ at $x=4,\ldots,13$.

\medskip

\noindent\textit{GF, column 13:}\ \ $-1$ under index $x=13$.

\noindent\textit{First factor (column 13):}\ \ row $t=1$: $1$ at $x=13$.

\noindent\textit{Second factor (column 13):}\ \ row $t=1$: $1,1,2,4,8,17,39,89,211,507,1238,3056$ at $x=2,\ldots,13$;\ \ row $t=2$: $1,1,3,6,15,33,82,194,482,1188,2988$ at $x=3,\ldots,13$;\ \ row $t=3$: $1,1,3,7,17,40,102,249,634,1594$ at $x=4,\ldots,13$.

\bigskip
\noindent\textbf{[p.~459]}\quad
I annex the following two Tables of (centre- and bicentre-) trees as far as I have completed them.

\medskip
\begin{center}
\textbf{Table A.}\\
\textit{Valency not greater than}
\end{center}

\begin{center}
\renewcommand{\arraystretch}{1.1}
\begin{tabular}{c|ccccccccc|c}
\hline
Knots & 0 & 1 & 2 & 3 & 4 & 5 & 6 & 7 & 8 & Gen. \\
 & & & Oxygen. & Boron. & Carbon. & & & & & \\
\hline
1  & 1 & 1 & 1 & 1 & 1 & 1 & 1 & 1 & 1 & 1 \\
2  &   & 1 & 1 & 1 & 1 & 1 & 1 & 1 & 1 & 1 \\
3  &   &   & 1 & 1 & 1 & 1 & 1 & 1 & 1 & 1 \\
4  &   &   & 1 & 2 & 2 & 2 & 2 & 2 & 2 & 2 \\
5  &   &   & 1 & 2 & 3 & 3 & 3 & 3 & 3 & 3 \\
6  &   &   & 1 & 4 & 5 & 6 & 6 & 6 & 6 & 6 \\
7  &   &   & 1 & 6 & 9 & 10 & 11 & 11 & 11 & 11 \\
8  &   & 1 & 1 & 11 & 18 & 21 & 22 & 23 & 23 & 23 \\
9  &   &   & 1 & 18 & 35 & 42 & 45 & 46 & 47 & 47 \\
10 &   &   & 1 & 37 & 75 &   &   &   &   & 106 \\
11 &   &   & 1 & 66 & 159 &   &   &   &   & 235 \\
12 &   &   & 1 & 135 & 367 &   &   &   &   & 551 \\
13 &   &   & 1 & 265 & 799 &   &   &   &   & 1301 \\
\hline
\end{tabular}
\end{center}

\medskip
\begin{center}
\textbf{Table B.}\\
\textit{Actual Valency.}
\end{center}

\begin{center}
\renewcommand{\arraystretch}{1.1}
\begin{tabular}{c|ccccccccc}
\hline
Knots & 0 & 1 & 2 & 3 & 4 & 5 & 6 & 7 & 8 \\
\hline
1  & 1 &   &   &   &   &   &   &   &   \\
2  &   & 1 &   &   &   &   &   &   &   \\
3  &   &   & 1 &   &   &   &   &   &   \\
4  &   &   & 1 & 1 &   &   &   &   &   \\
5  &   &   & 1 & 1 & 1 &   &   &   &   \\
6  &   &   & 1 & 3 & 1 & 1 &   &   &   \\
7  &   &   & 1 & 5 & 3 & 1 & 1 &   &   \\
8  &   &   & 1 & 10 & 7 & 3 & 1 & 1 &   \\
9  &   &   & 1 & 17 & 17 & 7 & 3 & 1 & 1 \\
10 &   &   & 1 & 36 & 38 &   &   &   &   \\
11 &   &   & 1 & 65 & 93 &   &   &   &   \\
12 &   &   & 1 & 134 & 232 &   &   &   &   \\
13 &   &   & 1 & 264 & 534 &   &   &   &   \\
\hline
\end{tabular}
\end{center}

\medskip
In A, the columns 2, 3, 4, and the last column are the totals given by the Tables IV., VI., VIII., and II., and the remaining numbers of columns 5, 6, 7, 8 have been found by trial; and, in B, the several columns are the differences of the

\bigskip
\noindent\textbf{[p.~460]}\quad
columns of A. The signification is obvious; for instance, if the number of knots is $=9$, then Table A, if the valency, or the maximum number of branches from a knot, is $=2,\,3,\,4,\,5,\,6,\,7,\,8$ or any greater number,
\[
\text{No.\ of trees }=\;1,\;18,\;35,\;42,\;45,\;46,\;47:
\]
viz.\ with $9$ knots the tree can have at most $8$ branches from a knot, so that the number of trees having at most $8$ branches from a knot is $=47$, the whole number of trees with $9$ knots; and so the number of knots being as before $=9$, Table B shows that the number of $47$ is made up of the numbers
\[
1,\;17,\;17,\;7,\;3,\;1,\;1;
\]
viz.\ $1$ is the No.\ of trees, at most $2$ branches from a knot,
\begin{align*}
17 &\quad\text{trees}\quad\text{at least one $3$-branch knot,}\\
17 &\quad\text{trees}\quad\text{,, ,, ,, $4$-branch knot,}\\
7  &\quad\text{trees}\quad\text{,, ,, ,, $5$-branch knot,}\\
3  &\quad\text{trees}\quad\text{,, ,, ,, $6$-branch knot,}\\
1  &\quad\text{tree}\quad\text{,, ,, ,, $7$-branch knot,}\\
1  &\quad\text{tree}\quad\text{,, ,, ,, $8$-branch knot,}
\end{align*}
and we have accordingly
\[
1+1+2+3+6+11+23+47.
\]

I annex also a plate showing the figures of the trees of $1,\,2,\,3,\ldots,\,9$ knots, classified according to their altitudes and number of main branches; and as to the bicentre-trees, according to the number of main branches from each point of the bicentre. The affixed numbers show in each case the greatest number of branches from a knot; so that when this is $(2)$, the knots may be oxygen-, boron-, carbon-, \&c., atoms; when $(3)$, boron-, carbon-, \&c., atoms; when $(4)$, carbon-, \&c., atoms; and so on.

\begin{center}
\par\noindent\fbox{\parbox{0.9\textwidth}{\footnotesize [Plate — see scan: of trees follows in the original printing; not reproduced here.]}}\par
\end{center}

\bigskip
\hrule
\bigskip

%% ========================================================================
%% PAPER [611] -- REPORT OF THE COMMITTEE ON MATHEMATICAL TABLES
%%                Printed pp. 461--478 (PDF pp. 481--500, with pp. 481-482 blank)
%% ========================================================================
\section*{\textbf{611.} Report of the Committee on Mathematical Tables}
\addcontentsline{toc}{section}{\textbf{611.} Report of the Committee on Mathematical Tables (pp.~461--478)}

\textit{Consisting of Professor Cayley, F.R.S., Professor Stokes, F.R.S., Professor Sir W.~Thomson, F.R.S., Professor H.~J.~S.~Smith, F.R.S., and J.~W.~L.~Glaisher, F.R.S.}

\medskip
\noindent\textit{[From the Report of the British Association for the Advancement of Science (1875), pp.~305--336.]}

\bigskip
\noindent\textbf{[p.~461]}\quad
The present Report (say Report 1875) is in continuation of that by Mr Glaisher, published in the volume for 1873, and here cited as Report 1873.

Report 1873 extends to all those tables which are at p.~3 (\textit{l.c.}) included under the headings:---
\begin{itemize}
\item[] A,\ auxiliary for non-logarithmic calculation, 1, 2, 3;
\item[] B,\ logarithmic and circular, 4, 5, 6;
\item[] C,\ exponential, 7, 8 (but only partially to C.~8), other than those tables of C referred to as ``h$\,.\,$l$\tan(45^\circ+\tfrac{1}{2}\phi)$''; and also partially (see Art.~24, pp.~81--83) to the tables included under the heading ``E.~11, transcendental constants $e,\,\pi,\,\gamma,$ \&c., and their powers and functions.''
\end{itemize}

A future Report will comprise the tables, or further tables, included under the headings:---
\begin{itemize}
\item[] C.~8.\ Hyperbolic antilogarithms $(e^x)$ and $h.\,l\tan(45^\circ+\tfrac{1}{2}\phi)$, and hyperbolic sines, cosines, \&c.
\item[] D.\ Algebraic constants.
\item[\phantom{D.}] \ \ 9.\ Accurate integer or fractional values.\ \ Bernoulli's Numbers, $\Delta^n 0^m$, \&c.\ \ Binomial coefficients.
\item[\phantom{D.}] 10.\ Decimal values auxiliary to the calculation of series.
\item[] E.\ 11.\ Transcendental constants $e,\,\pi,\,\gamma,$ \&c., and their powers and functions.
\end{itemize}

\bigskip
\noindent\textbf{[p.~462]}\quad
The present Report (1875) comprises the tables included under the headings:---
\begin{itemize}
\item[] F.\ Arithmological.
\item[\phantom{F.}] 12.\ Divisors and prime numbers.\ \ Prime roots.\ \ The Canon arithmeticus, \&c.
\item[\phantom{F.}] 13.\ The Pellian equation.
\item[\phantom{F.}] 14.\ Partitions.
\item[\phantom{F.}] 15.\ Quadratic forms $a^2+b^2$, \&c., and partitions of numbers into squares, cubes, and biquadrates.
\item[\phantom{F.}] 16.\ Binary, ternary, \&c., quadratic and higher forms.
\item[\phantom{F.}] 17.\ Complex theories;
\end{itemize}
which divisions are herein referred to, for instance, as [F.~12.\ Divisors, \&c.].

Report 1873 consists of six sections (\S) divided into articles, which are separately numbered (see contents, p.~174); the present Report 1875 forms a single section (\S~7), divided in like manner into articles, which are separately numbered; but besides this the paragraphs are numbered, and that continuously, through the present Report 1875, so that any paragraph may be cited as Report 1875, No.~---, as the case may be.

\bigskip
\begin{center}
\textbf{[F.~12.\ Divisors, \&c.]\ \ Divisors and Prime Numbers.\ \ Art.~I.}
\end{center}

\textbf{1.}\ \ As to divisors and prime numbers see Report 1873, Art.~8 (Tables of Divisors---factor tables---and Tables of Primes), pp.~34--40.   The tables there referred to, such as Chernac, Burckhardt, Dase, Dase and Rosenberg, are chiefly tables running up to very high numbers (the last of them the ninth million): wherein, to save space, multiples of $2,\,3,\,5$ are frequently omitted, and in some of them only the least divisor is given.   It would be for many purposes convenient to have a small table, going up say to 10,000, showing in every case all the prime factors of the number.   Such a table might be arranged, $500$ numbers in a page, in some such form as the following:

\begin{center}
\textit{Factor Table 1 to 500}
\end{center}

\begin{center}
\renewcommand{\arraystretch}{1.1}
\begin{tabular}{c|ccccccccccc}
\hline
 & 0 & 1 & 2 & 3 & 4 & 5 & 6 & 7 & 8 & 9 \\
\hline
39 & $2.3.5.13$ & $17.23$ & $2^2.7^2$ & $3.131$ & $2.197$ & $5.79$ & $2^2.3^2.11$ & $397^{*}$ & $2.199$ & $3.7.19$ \\
\hline
\end{tabular}
\end{center}

\noindent where the top line shows the units, and the left-hand column the remaining figures, viz.\ the specimen exhibits the composition of the several numbers from $390$ to $399$: a prime number, e.g.\ $397$, would be sufficiently indicated by the absence of any decomposition, or it may be further indicated by an asterisk.

It may be noticed that, in the theory of numbers, the decomposition is specially required when the next following number is a prime, viz.\ that of $p-1$, $p$ being a

\bigskip
\noindent\textbf{[p.~463]}\quad
prime: also, that this is given incidentally, for prime numbers $p$ up to $1000$, in Jacobi's \textit{Canon Arithmeticus}, \textit{post}, No.~20, and up to $15{,}000$ in Reuschle's Tables, V.~(a, b, c) \textit{post}, No.~22.

\textbf{2.}\ \ It may be proper to remark here that any table of a binary form is really a factor-table in the complex theory connected with such binary form.   Thus in a table of the form $a^2+b^2$, a number of this form has a factor $a+bi$ ($i=\sqrt{-1}$ as usual); and the table, in fact, shows the complex factor $a+bi$ of the number in question: a well arranged table would give all the prime complex factors $a+bi$ of the number.   But as to this more hereafter; at present, we are concerned with the real theory only, not with any complex theory.

\textbf{3.}\ \ Connected with a factor-table, we have (i) a Table of the number of less relative primes; viz.\ such a table would show, for every number, the number of inferior integers having no common factor with the number itself.   The formula is a well-known one: for a number $N=a^\alpha b^\beta c^\gamma\ldots$, ($a,\,b,\ldots$ the distinct prime factors of $N$), the number of less relative primes is
\[
\totient(N),\ =a^{\alpha-1}b^{\beta-1}\cdots(a-1)(b-1)\cdots,
\]
or, what is the same thing, $=N\!\left(1-\dfrac{1}{a}\right)\!\left(1-\dfrac{1}{b}\right)\cdots$\ \ A small table $(N=1\text{ to }100)$, occupying half a page, is given by

Euler, \textit{Op.\ Arith.\ Coll.}\ t.~ii.\ p.~128; viz.\ this is $\totient 1=0,\;\totient 2=1,\ldots,\;\totient 100=40$.

\textbf{4.}\ \ But it would be interesting to have such a table of the same extent with the proposed factor-table.   The table might be of like form; for instance,

\begin{center}
\textit{Number of less relative Primes Table 1 to 500}
\end{center}

\begin{center}
\renewcommand{\arraystretch}{1.1}
\begin{tabular}{c|ccccccccccc}
\hline
 & 0 & 1 & 2 & 3 & 4 & 5 & 6 & 7 & 8 & 9 \\
\hline
39 & $112$ & $192$ & $144$ & $292$ & $84$ & $232$ & $144$ & $198$ & $148$ & $264$ \\
\hline
\end{tabular}
\end{center}

It would be of still greater interest to have an inverse table showing the values of $N$ which belong to a given value of $\totient(N)$; for instance,
\[
\begin{array}{r|l}
\totient= & N= \\\hline
40 & 41,\;55,\;75,\;82,\;88,\;100,\;110, \\
42 & 43,\;49,\;86,\;98, \\
44 & 69,\;92, \\
46 & 47,\;94, \\
48 & 65,\;104,\;105,\;112,\\
\vdots & \vdots
\end{array}
\]
where, observe, that $\totient$ is of necessity even.

\bigskip
\noindent\textbf{[p.~464]}\quad
\textbf{5.}\ \ Again, connected with a factor-table, we have (ii) a Table of the Sum of the divisors of a Number.   The formula is also a well-known one; for a number $N=a^\alpha b^\beta\ldots$, ($a,\,b,\ldots$ the distinct prime factors of $N$), the required sum
\[
\sumdiv N\;=\;(1+a+\cdots+a^\alpha)(1+b+\cdots+b^\beta)\cdots,
\]
or, what is the same thing,
\[
=\frac{a^{\alpha+1}-1}{a-1}\cdot\frac{b^{\beta+1}-1}{b-1}\cdots,
\]
where, observe, that the number itself is reckoned as a divisor.

\textbf{6.}\ \ Such a table was required by Euler in his researches on Amicable Numbers (see \textit{post}, No.~10), and he accordingly gives one of a considerable extent, viz.

Euler, \textit{Op.\ Arith.\ Coll.}\ t.~i.\ pp.~104--109.

It is to be remarked that, inasmuch as $\sumdiv N$ is obviously $\sumdiv a^\alpha\cdot\sumdiv b^\beta\ldots$, the function need only be tabulated for the different integer powers $a^\alpha$ of each prime number $a$.   The range of Euler's table is as follows:
\[
\begin{array}{rl}
a= & a= \\
2  & \text{to }36, \\
3  & \;\,,, \;\,15, \\
5  & \;\,,, \;\,\phantom{0}9, \\
7  & \;\,,, \;\,10, \\
11 & \;\,,, \;\,\phantom{0}9, \\
13 & \;\,,, \;\,\phantom{0}7, \\
17 & \;\,,, \;\,\phantom{0}5, \\
19 & \;\,,, \;\,\phantom{0}5, \\
23 & \;\,,, \;\,\phantom{0}4, \\
29\text{ to }997 & \;\,,, \;\,\phantom{0}3,
\end{array}
\]
viz.\ for the several prime numbers from $29$ to $997$ the table gives $\sumdiv a$, $\sumdiv a^2$ and $\sumdiv a^3$.

And it is to be noticed that the values of the sum are exhibited, decomposed into their prime factors: thus a specimen of the table is
\[
\begin{array}{c|c}
\textit{Num.} & \textit{Summa Divisorum.}\\\hline
139    & 2^2.5.7 \\
139^2  & 3.13.499 \\
139^3  & 2^3.5.7.9661
\end{array}
\]

\bigskip
\noindent\textbf{[p.~465]}\quad
\textbf{7.}\ \ The form of the above table is adapted to its particular purpose (the theory of amicable numbers); but Euler gives also,

Euler, \textit{Op.\ Arith.\ Coll.}\ t.~i.\ p.~147---in the paper ``Observatio de Summis Divisorum,'' (1752), pp.~146--154,---a short table of about half a page, $N=1$ to $100$, of the form $\sumdiv 1=1$, $\sumdiv 2=3$, \ldots, $\sumdiv 100=217$.   The paper contains interesting analytical researches on the subject of $\sumdiv N$ which connect themselves with the theory of the Partition of Numbers.

\textbf{8.}\ \ It would be interesting to carry the last-mentioned table to the same extent as the proposed factor-table; and to add to it an inverse table, as suggested in regard to the number of less relative primes table.

\textbf{9.}\ \ \textit{Perfect Numbers.}\ ---A perfect number is a number which is equal to the sum of its divisors, the number itself not being reckoned as a divisor; e.g.\
\[
6=1+2+3,\quad\text{and}\quad 28=1+2+4+7+14.
\]
Such numbers are indicated by a table of the sums of divisors: $\sumdiv 6=12$, $\sumdiv 28=56$, these two being, as appears by the table, Art.~7, the only perfect numbers less than $100$.

\textbf{10.}\ \ \textit{Amicable Numbers.}\ ---These are pairs of numbers such that each is equal to the sum of the divisors of the other, not reckoning the other number as a divisor; that is, each has the same sum of divisors, the number being here reckoned as a divisor; say $\sumdiv A=B$, $\sumdiv B=A$; or, what is the same thing, $\sumdiv A=\sumdiv B\;(=A+B)$.   Thus for the numbers $220,\,284$,
\[
\begin{aligned}
\sumdiv 220 &= (1+2+4)(1+5)(1+11)-220, = 284, \\
\sumdiv 284 &= (1+2+4)(1+71)-284, = 220;
\end{aligned}
\]
or, what is the same thing,
\[
\sumdiv 220 = (1+2+4)(1+5)(1+11) = 504 = (1+2+4)(1+71) = \sumdiv 284.
\]

\textbf{11.}\ \ A catalogue of $61$ pairs of numbers is given by

Euler, \textit{Op.\ Arith.\ Coll.}\ t.~i.\ pp.~144--145; it occupies about one page.   The paper, ``De Numeris Amicabilibus,'' pp.~102--145, contains an elaborate investigation of the theory, by means whereof all but two of the pairs of numbers are obtained.   The first pair is the above-mentioned one, $2^2.5.11$ and $2^2.71$ ($=220$ and $284$); and the fifty-ninth pair is the high numbers
\[
3^2.7^2.13.19.53.6959\quad\text{and}\quad 3^2.7^2.13.19.179.2087.
\]

\bigskip
\noindent\textbf{[p.~466]}\quad
The last two pairs are referred to as ``form\ae\ diverse a precedentibus''; viz.\ these are
\[
\left\{\begin{array}{c}2^3.19.41 \\ 2^3.199\end{array}\right.\quad\text{and}\quad\left\{\begin{array}{c}2^3.41.467 \\ 2^3.19.233\end{array}\right.
\]

\textbf{12.}\ \ A Table of the Frequency of Primes is given by

Gauss, \textit{Tafel der Frequenz der Primzahlen}, \textit{Werke}, t.~ii.\ pp.~436--443; viz.\ this extends to $3{,}000{,}000$.

The first part, extending to $1{,}000{,}000$, $=1000$ thousand, shows how many primes there are in each thousand: a specimen is
\[
\begin{array}{rl}
1, & 168;\\
2, & 135;\\
3, & 127;\\
4, & 120;\\
5, & 119;\\
\&c. &
\end{array}
\]
viz.\ in the first thousand there are $168$ primes, in the second thousand $135$ primes, and so on.

For the second and third millions the frequency is given for each ten thousand: a specimen is
\begin{center}
\textit{1,000,000 to 1,100,000.}
\end{center}

\begin{center}
\renewcommand{\arraystretch}{1.0}
\begin{tabular}{r|rrrrrrrrrr|r}
\hline
 & 0 & 1 & 2 & 3 & 4 & 5 & 6 & 7 & 8 & 9 &  \\\hline
1  & 1 &    &    &    &    &    &    &    &    &    & 1 \\
2  &   & 1  &    &    & 1  &    & 1  & 1  &    &    & 4 \\
3  & 1 & 4  & 2  & 2  & 3  & 1  & 2  & 3  & 3  &    & 21 \\
4  & 2 & 8  & 5  & 4  & 3  & 6  & 9  & 4  & 5  & 8  & 54 \\
5  & 11& 10 & 8  & 18 & 12 & 10 & 10 & 12 & 15 & 8  & 114 \\
6  & 14& 14 & 18 & 21 & 16 & 22 & 19 & 15 & 17 & 15 & 171 \\
7  & 26& 17 & 23 & 23 & 24 & 24 & 17 & 22 & 20 & 21 & 217 \\
8  & 19& 19 & 21 & 7  & 14 & 15 & 20 & 17 & 15 & 17 & 164 \\
9  & 11& 13 & 9  & 13 & 14 & 14 & 12 & 13 & 11 & 16 & 126 \\
10 & 8 & 6  & 8  &    & 9  & 5  & 9  & 9  & 7  & 9  & 71 \\
11 & 6 & 6  & 4  & 6  & 3  & 1  & 3  & 1  & 4  & 5  & 39 \\
12 & 1 & 1  & 2  & 1  & 1  & 1  & 2  & 2  & 1  &    & 12 \\
13 & 1 & 1  &    &    & 1  &    & 1  & 1  & 1  &    & 6 \\
14 &   &    &    &    &    &    &    &    &    &    &   \\
15 &   &    &    &    &    &    &    &    &    &    &   \\
16 &   &    &    &    &    &    &    &    &    &    &   \\\hline
   & 752 & 719 & 732 & 700 & 731 & 698 & 713 & 722 & 706 & 737 & 7210 \\\hline
\end{tabular}
\end{center}
\[
\int\frac{dx}{\log x} = 7212\cdot 99\,;
\]

\bigskip
\noindent\textbf{[p.~467]}\quad
viz.\ in the interval $1{,}000{,}000$ to $1{,}010{,}000$, $100$ hundreds, there is $1$ hundred containing $1$ prime, there are $2$ hundreds each containing $4$ primes, $11$ hundreds each containing $5$ primes,\,\ldots, $1$ hundred containing $13$ primes, so that, as
\[
\begin{aligned}
1\times 1   &=  1,\\
4\times 2   &=  8,\\
5\times 11  &=  55,\\
&\;\vdots\\
13\times 1  &=  13,\\\hline
100\phantom{XXX} &\phantom{=\;}\,752,
\end{aligned}
\]
the whole $10{,}000$ contains $752$ primes; the next $10{,}000$ contains $719$ primes, and so on; the whole $100{,}000$ thus containing $752+719+\cdots\;=7210$ primes, which number is at the foot compared with the theoretic approximate value
\[
\int\frac{dx}{\log x}\quad(\text{limits }1{,}000{,}000\text{ to }1{,}010{,}000) = 7212\cdot 99.
\]
The integral in question is represented by the notation $\mathrm{Li}.\,x$ or $\mathrm{li}.\,x$.

p.~443.\ \ We have the like tables $1{,}000{,}000$ to $2{,}000{,}000$ and $2{,}000{,}000$ to $3{,}000{,}000$, showing for each $100{,}000$ how many hundreds there are containing $0$ prime, $1$ prime, $2$ primes, up to (the largest number) $17$ primes.

\textbf{13.}\ \ It is noticed that
\begin{quote}
the $26{,}379$th hundred contains no prime,\\
the $27{,}050$th hundred contains $17$ primes.
\end{quote}

It may be observed that, if $N=2.3.5\ldots p$, the product of all the primes up to $p$, then each of the numbers $N+1$ and $N+q$ (if $q$ be the prime next succeeding $p$) is or is not a prime; but the intermediate numbers $N+2,\,N+3,\ldots,\,N+q-1$ are certainly composite; viz.\ we thus have at least $q-2$ consecutive composites.   To obtain in this manner $99$ consecutive composites, the value of $N$ would be $=2.3.5\ldots 97$, viz.\ this is a number far exceeding $2{,}637{,}900$; but, in fact, the hundred numbers $2{,}637{,}901$ to $2{,}638{,}000$ are all composite.

Legendre, in his \textit{Essai sur la Th\'eorie des Nombres} (1st edit., 1798; 2nd edit., 1808, supplement, 1816: references to this edition), gives for the number of primes inferior to a given limit $x$ the approximate formula
\[
\frac{x}{\log x-1\cdot 08366},
\]
and p.~394, and supplement, p.~62, he compares for each $10{,}000$ up to $100{,}000$, and for each $100{,}000$ up to $1{,}000{,}000$, the values as computed by this formula with the actual numbers of primes exhibited by the tables of Vega and Chernac.   Thus for $x=1{,}000{,}000$, the computed value is $78{,}543$, the actual value $78{,}493$.

\bigskip
\noindent\textbf{[p.~468]}\quad
He shows, p.~414, that the number of integers, which are less than $n$ and are not divisible by any of the numbers $\theta,\,\lambda,\,\mu,\ldots$, is approximately
\[
=n\!\left(1-\tfrac{1}{\theta}\right)\!\left(1-\tfrac{1}{\lambda}\right)\!\left(1-\tfrac{1}{\mu}\right)\cdots,
\]
and taking $\theta,\,\lambda,\,\mu,\ldots$ the successive primes $3,\,5,\,7,\ldots$ he gives the values of the function in question, or, say, the function
\[
\frac{2}{3}\cdot\frac{4}{5}\cdot\frac{6}{7}\cdot\frac{10}{11}\cdots\frac{\omega-1}{\omega},
\]
$\omega$ a prime, for the several prime values $\omega=3$ to $1229$ in the Table IX.\ (one page) at the end of the work.

\textbf{14.}\ \ A table of frequency is given by

Glaisher, J.~W.~L., \textit{British Association Report} for $1872$, p.~20.   This gives for the second and the ninth millions, respectively divided into intervals of $50{,}000$, the actual number of primes in each interval, as compared with the theoretic value $\mathrm{li}\,x'-\mathrm{li}\,x$; and also deduced therefrom, by the formula $\log\tfrac{1}{2}(x'+x)$, a table of the average interval between two consecutive primes; this average interval increases very slowly: at the beginning and the end of the second million the values are $13\cdot 76$ and $14\cdot 58$ (theoretic values $13\cdot 84$ and $14\cdot 50$); at the beginning and the end of the ninth million $16\cdot 02$ and $15\cdot 95$ (theoretic values $15\cdot 90$ and $16\cdot 01$).

\textbf{15.}\ \ Coming under the head of Divisor Tables, some tables by Reuschle and Gauss may be here referred to.   These are:

Reuschle, \textit{Mathematische Abhandlung, zahlentheoretische Tabellen sammt einer dieselben treffenden Correspondenz mit der verewigten C.~G.~J.~Jacobi}, $4^{\circ}$, pp.~1--61\footnote{Titlepage missing in my copy; but I find from Prof.\ Kummer's notice of the work, Crelle, t.~liii.\ (1857), p.~379, that it appeared as a Programm of the Stuttgart Gymnasium, Michaelmas, 1856, and was separately printed by Liesching and Co., Stuttgart.} (1856).   The tables belonging to the present subject are

A.\ Tafeln zur Zerlegung von $a^n-1$ (pp.~18--22).

I.\ Table of the prime factors of $10^n-1$, viz.

(a.\ pp.~18--19.)\ \ Complete decomposition of $10^n-1$, $n=1$ to $42$: and $10^n+1$, $n=1$ to $21$.   Some values of $n$ are omitted.

A specimen is
\[
\begin{aligned}
10^{14}-1 &= 3^2.53.79.265371653, \\
10^{12}+1 &= 11.139.10\,5831\,3049.
\end{aligned}
\]

(b.\ p.~19.)\ \ List of the specific prime factors $f$ of $10^n-1$, or the prime factors of the residue after separation of the analytical factors, for those values of $n$ for which the complete decomposition is unknown, and omitting those values for which no factor is known, $n=25$ to $243$.

\bigskip
\noindent\textbf{[p.~469]}\quad
A specimen is
\[
\begin{array}{c|c}
n & f \\\hline
25 & 21401
\end{array}
\]
The meaning seems to be, residue of $10^{25}-1$ is $1+10^5+10^{10}+10^{15}+10^{20}$ and this contains the prime factor $21401$; but it is not clear why this is the ``specific prime factor.''

II.\ Prime factors of $a^n-1$ for different values of $a$ and $n$.

(a.\ p.~20) gives for $41$ values of $a$ ($2,\,3$, \&c.\ at intervals to $100$) and for the following values of $n$ the decompositions of the residues or specific factors of $a^n-1$, viz.\ these are
\[
\begin{aligned}
n &= 1,\ a-1 \\
  &\phantom{=\;}2,\ a+1 \\
  &\phantom{=\;}3,\ a^2+a+1 \\
  &\phantom{=\;}6,\ a^2-a+1 \\
  &\phantom{=\;}4,\ a^2+1 \\
  &\phantom{=\;}5,\ a^4+a^3+a^2+a+1 \\
  &\phantom{=\;}10,\ a^4-a^3+a^2-a+1 \\
  &\phantom{=\;}8,\ a^4+1 \\
  &\phantom{=\;}12,\ a^4-a^2+1.
\end{aligned}
\]

A specimen is
\[
\begin{array}{c|c|c|c|c|c|c}
   & a-1 & a^2-1 & a^3-1 & a^6-1 & a^4-1 & a^5-1 \\\hline
10 & 3^2 & 11 & 3^2.37 & 7.13 & 101 & 41.271 \\
   & & & a^{10}-1 &  &  & \\
   & & & 9091 & & 73.137 & 9901
\end{array}
\]

(b.\ p.~21.)\ \ Specific prime factors for the numbers $2,\,3,\,5,\,6,\,7,\,10$, (the powers $4,\,8,\,9$ being omitted as coming under $2$ and $3$), for the exponents $1$ to $42$.

A specimen is
\[
\begin{array}{c|c|c|c|c|c|c}
   & 2^n-1 & 3^n-1 & 5^n-1 & 6^n-1 & 7^n-1 & 10^n-1 \\\hline
19 & 524287 & 1597.363889 & 191.x & 191.x & 419.x & --
\end{array}
\]
where the $x$ denotes that the other factor is not known to be prime.   And so, where no number is given, as in $10^{19}-1$, it is not known whether the number $(=1+10+10^2+\cdots+10^{18})$ is or is not prime.

Addition, p.~22.\ \ For $a=2$, the complete decomposition of the prime factor of $2^n-1$ is given for values of $n,\,=44,\,45,\ldots$ at intervals to $156$.

A specimen is
\[
\begin{array}{c|c}
n & f \\\hline
44 & 397.2113
\end{array}
\]
viz.
\[
2^{22}-2^{11}-\cdots-2^2+1\;=\;838861\;=\;397.2113.
\]

$n=31$, Fermat's prime.\ \ $n=37$, the first case for which the decomposition is not given completely.   $n=41$, the first case for which no factor is known.

\bigskip
\noindent\textbf{[p.~470]}\quad
\textbf{16.}\ \ Gauss, \textit{Tafel zur Cyclotechnie}, \textit{Werke}, t.~ii.\ pp.~478--495, shows, for $2452$ numbers of the several forms $a^2+1,\,a^2+4,\,a^2+9,\ldots,\,a^2+81$, the values of $a$ such that the number in question is a product of prime factors no one of which exceeds $200$, and exhibits all the odd prime factors of each such number.   The table is in nine parts, \textit{zerlegbare $a^2+1$}, \textit{zerlegbare $a^2+4$}, \&c., with to each part a subsidiary table, as presently mentioned.   Thus a specimen is
\[
\begin{array}{c|l}
\multicolumn{2}{c}{\text{\textit{zerlegbare} }a^2+9}\\\hline
1 & 5\\
2 & 13\\
4 & 5.5\\
5 & 17\\
7 & 29\\
8 & 73\\
\vdots & \\
1411168679 & 5.5.13.17.17.89.113.157.173.197.197
\end{array}
\]
viz.
\[
\begin{aligned}
1^2+9, &\text{ odd prime factor is }5,\\
2^2+9, &\quad\quad\,,, \quad\,,, \quad\quad 13,\\
4^2+9, &\quad\quad\,,, \quad\,,, \text{ factors are }5,5,
\end{aligned}
\]
and so on.

And the subsidiary table is
\[
\begin{array}{c|l}
 5 & 1,\,4,\,79 \\
13 & 2,\,11,\,41 \\
17 & 5,\,29,\,46,\,379,\,1042
\end{array}
\]
showing that the numbers $a$ for which the largest factor is $5$ are $1,\,4,\,79$; those for which it is $13$ are $2,\,11,\,41$; and so on.

The object of the table is explained in the \textit{Bemerkungen}, (\textit{l.c.}, p.~523), by Schering, the editor of the volume, viz.\ it is to facilitate the calculation of the circular arcs the cotangents of which are rational numbers.   To take a simple example, it appears to be by means of it that Gauss obtained, among other formul\ae, the following:
\[
\tfrac{\pi}{4} = 12\,\arctan\tfrac{1}{38} + 8\,\arctan\tfrac{1}{57} - 5\,\arctan\tfrac{1}{239},
\]
and
\[
\tfrac{\pi}{4} = 12\,\arctan\tfrac{1}{18} + 20\,\arctan\tfrac{1}{57} + 7\,\arctan\tfrac{1}{239} + 24\,\arctan\tfrac{1}{268}.
\]

\bigskip
\noindent\textbf{[p.~471]}\quad
\begin{center}
[F.~12.\ Divisors, etc.] continued.\ \ Prime Roots.\ \ The Canon Arithmeticus, Quadratic residues.\ \ Art.~II.
\end{center}

\textbf{17.}\ \ \textit{Prime Roots.}\ ---Let $p$ be a prime number; then there exist $\totient(p-1)$ inferior integers $g$, such that all the numbers $1,\,2,\ldots,\,p-1$ are, to the modulus $p$,
\[
=1,\,g,\,g^2,\,\ldots,\,g^{p-2}\quad(g^{p-1}\text{ is of course }=1).
\]
This being so, $g$ is said to be a \textit{prime root} of $p$; and moreover the several numbers $g^\alpha$, where $\alpha$ is any number whatever less than and prime to $p-1$, constitute the series of the $\totient(p-1)$ prime roots of $p$.   It may be added that, if $\beta$ be an integer number less than $p-1$, and having with it a greatest common measure $=k$, so that $\beta=\dfrac{p-1}{k}\cdot t$, then
\[
(g^\beta)^{\frac{p-1}{k}} = g^{t(p-1)} = 1,\;(\text{since $t$ is an integer, and }g^{p-1}=1)
\]
then $g^\beta$ has the indicatrix $\dfrac{p-1}{k}$: the prime roots are those numbers which have the indicatrix $p-1$.

The like theory exists as to any number $N$ of the form $p^m$ or $2p^m$.   There are here $\totient(N),\,=N\!\left(1-\dfrac{1}{p}\right)$ or $\tfrac{1}{2}N\!\left(1-\dfrac{1}{p}\right)$, in the two cases respectively, numbers less than $N$ and prime to it; and we have then $\totient\!\left(\totient(N)\right)$ numbers $g$ such that, to the modulus $N$, all these numbers are $=1,\,g,\,g^2,\,\ldots,\,g^{\totient(N)-1}$ ($g^{\totient(N)}$ is of course $=1$).   This being so, $g$ may be regarded as a prime root of $N$ ($=p^m$ or $2p^m$, as the case may be); and moreover the several numbers $g^\alpha$, where $\alpha$ is any number whatever less than and prime to $\totient(N)$, constitute the series of the $\totient(\totient(N))$ prime roots of $N$.   Thus $N=3^2=9$, $\totient(N)=6$; we have
\[
1,\,2^1,\,2^2,\,2^3,\,2^4,\,2^5,
\]
\[
=1,\,2,\,4,\,8,\,7,\,5,\ \text{mod.\ 9};
\]
or the prime roots of $9$ are $2^1$ and $2^5$, $=2$ and $5$.

So also $N=2.3^2=18$, $\totient(N)=6$; we have
\[
1,\,5^1,\,5^2,\,5^3,\,5^4,\,5^5,
\]
\[
=1,\,5,\,7,\,17,\,13,\,11,\ \text{mod.\ 18};
\]
and $5^1$ and $5^5$, $=5$ and $11$ are the prime roots of $18$.

\textbf{18.}\ \ A small table of prime roots, $p=3$ to $37$, is given by

Euler, \textit{Op.\ Arith.\ Coll.}\ t.~i.\ pp.~525--526.   The Memoir is entitled ``Demonstrationes circa residua e divisione potestatum per numeros primos resultantia,'' pp.~516--537 (1772).

\bigskip
\noindent\textbf{[p.~472]}\quad
\textbf{19.}\ \ A table, $p$ and $p^m$, $3$ to $97$, is given by

Gauss, ``Disquisitiones Arithmeticae,'' 1801, (\textit{Werke}, t.~i.\ p.~468).   This gives in each case a prime root, and shows the exponents in regard thereto of the several prime numbers less than $p$ or $p^m$.   Thus a specimen is
\[
\begin{array}{c|c|cccccccccc}
  &   & 2 & 3 & 5 & 7 & 11 & 13 & 17 & 19 & 23 & 29\ \&c. \\\hline
27 & 2 & 1 & * & 5 & 16 & 13 & 8 & 15 & 12 & 11 & \\
29 & 10& 11& 27 & 18 & 20 & 23 & 2 & 7 & 15 & 24 &
\end{array}
\]
viz.\ for $27$ we have $2$ a prime root, and $2=2^1$, $5=2^5$, $7=2^{16}$, $11=2^{13}$, \&c.; and so also for $29$ we have $10$ a prime root, and $2=10^{11}$, $3=10^{27}$, $5=10^{18}$, \&c.

\textbf{20.}\ \ Small tables are probably to be found in many other places; but the most extensive and convenient table is Jacobi's \textit{Canon Arithmeticus}, the complete title of which is

\textit{Canon Arithmeticus sive tabul\ae\ quibus exhibentur pro singulis numeris primis vel primorum potestatibus infra 1000 numeri ad datos indices et indices ad datos numeros pertinentes.   Edidit C.\ G.\ J.\ Jacobi.   Berolini, 1839.\ \ $4^\circ$.}

The contents are as follows:
\[
\begin{array}{lr}
\textit{Introductio} & i\text{ to }\text{xl} \\
\textit{Tabul\ae\ numerorum ad indices datos pertinentium et indicum} & \\
\quad\textit{numero dato correspondentium pro modulis primis minoribus} & \\
\quad\textit{quam 1000} & 1\text{--}221 \\
\textit{Tabul\ae\ residuorum et indicum sibi mutuo respondentium pro} & \\
\quad\textit{modulis minoribus quam 1000 qui sunt numerorum primorum} & \\
\quad\textit{potestates} & 222\text{--}238 \\
\textit{Hujus tabulae ea pars quae pertinet ad modulos form\ae\ $2^n$, invenitur} & 239\text{--}240.
\end{array}
\]

The following is a specimen of the principal tables:
\[
p=19,\ p-1=2.3^2.
\]

\textit{Numeri.}
\[
\begin{array}{c|ccccccccc}
I & 1 & 2 & 3 & 4 & 5 & 6 & 7 & 8 & 9 \\\hline
  & 10& 5 & 12& 6 & 3 & 11& 15& 17& 18 \\
1 & 9 & 14& 7 & 13& 16& 8 & 4 & 2 & 1
\end{array}
\]

\textit{Indices.}
\[
\begin{array}{c|ccccccccc}
N & 1 & 2 & 3 & 4 & 5 & 6 & 7 & 8 & 9 \\\hline
  & 18& 17& 5 & 16& 2 & 4 & 12& 15& 10 \\
1 & 1 & 6 & 3 & 13& 11& 7 & 14& 8 & 9
\end{array}
\]

\bigskip
\noindent\textbf{[p.~473]}\quad
where the first table gives the values of the powers of the prime root $10$ (that $10$ is the root appears by its index being given as $=1$) to the modulus $19$, viz.\ $10^1=10$, $10^2=5$, $10^3=12$, \&c.; and the second table gives the index of the power to which the same prime root must be raised in order that it may be, to the modulus $19$, congruent with a given number: thus $10^{18}=1$, $10^{17}=2$, $10^5=3$, \&c.   The units of the index or number, as the case may be, are contained in the top line of the table, and the tens or hundreds and tens in the left-hand column.

\textbf{21.}\ \ There is given by

Jacobi, Crelle, t.~xxx.\ (1846), pp.~181, 182, a table of $m'$ for the argument $m$, such that
\[
1+g^m = g^{m'}\;(\text{mod.\ }p),\ p=7\text{ to }103,\text{ and }m=0\text{ to }102.
\]

A specimen is
\[
\begin{array}{c|ccccccccc|c}
p & 7 & 11 & 13 & 17 & 19 & 23 & 29 & 31 & 37 & \text{to }103 \\
g & 3 & 2  & 6  & 10 & 10 & 10 & 10 & 17 & 5  & \\\hline
m=11 & \cdot & \cdot & 6 & 4 & 7 & * & 27 & 21 & 34 &
\end{array}
\]
for instance, $p=19$, $1+10^{11}=10^7$ (mod.\ $19$).

Jacobi remarks that this table was calculated for him by his class during the winter course of 1836--37; and that, by means of the Canon Arithmeticus since published (in 1839), the same might easily be extended to all primes under $1000$.   In fact, for any such number $p$, putting any number of the table ``Indices'' $=m$, the next following number of the table gives the value of $m'$.

\textbf{22.}\ \ We have next, in Reuschle's Memoir (\textit{ante}, No.~15), the following relating to prime roots:

C.\ Tafeln f\"ur primitive Wurzeln und Hauptexponenten, oder V.~erweiterte und bereicherte Burkhardtsche Tafel, pp.~41--61, being divided into three parts; viz.\ these are

a.\ Table of the Hauptexponenten of the six roots $10,\,5,\,2,\,6,\,3,\,7$ for all prime numbers of the first $1000$, together with the least primitive root of each of these numbers (pp.~42--46).

A specimen is as follows:
\[
\begin{array}{c|c|cccccccccccc|c}
  &     & \multicolumn{2}{c}{10} & \multicolumn{2}{c}{5} & \multicolumn{2}{c}{2} & \multicolumn{2}{c}{6} & \multicolumn{2}{c}{3} & \multicolumn{2}{c|}{7} & w \\
p & p-1 & e & n & e & n & e & n & e & n & e & n & e & n &   \\\hline
53 & 2^2.13 & 13 & 4 & 52 & 1 & 52 & 1 & 26 & 2 & 52 & 1 & 26 & 2 & 2
\end{array}
\]
where $e$ is the Hauptexponent or indicatrix of the root ($10,\,5,\,2,\,6,\,3,\,7$, as the case may be), $n=\dfrac{p-1}{e}$, $w$ the least primitive root; thus
\[
p=53,\ 10^{13}=1,\ 5^{52}=1,\ 2^{52}=1,
\]

\bigskip
\noindent\textbf{[p.~474]}\quad
($2$ being accordingly the least prime root),
\[
6^{26}=1,\ 3^{52}=1,\ 7^{26}=1.
\]
The number $w$ of the last column is the least primitive root.   It is, of course, not always (as in the present case) one of the numbers $10,\,5,\,2,\,6,\,3,\,7$ to which the table relates: the first exception is $p=191$, $w=19$: the highest value of $w$ is $w=21$, corresponding to $p=409$.

b.\ The like table for the roots $10$ and $2$ for all prime numbers from $1000$ to $5000$, together with as convenient as possible a prime root (and in some cases two prime roots) for each such number (pp.~47--53).

A specimen is:
\[
\begin{array}{c|c|cccc|c}
  &     & \multicolumn{2}{c}{10} & \multicolumn{2}{c|}{2} &   \\
p & p-1 & e & n & e & n & w \\\hline
1289 & 2^3.7.23 & 92 & 14 & 161 & 8 & 6,\,11
\end{array}
\]
viz.\ here, mod.\ $1289$, $10^{92}=1$, $2^{161}=1$; and two prime roots are $6,\,11$.   We have thus by the present tables a prime root for every prime number not exceeding $5000$.

c.\ The like table for the root $10$ for all prime numbers between $5000$ and $15000$, (no column for $w$, nor any prime root given), pp.~53--61.

A specimen is
\[
\begin{array}{c|c|cc}
p & p-1 & e & n \\\hline
9859 & 2.3.31.53 & 3286 & 3
\end{array}
\]
viz., mod.\ $9859$, we have $10^{3286}=1$.   But in a large number of cases we have $n=1$, and therefore $10$ a prime root.   For example,
\[
9887\quad 2.4943\quad 9886\quad 1.
\]

\textbf{23.}\ \ For a composite number $n$, if $N=\totient(n)$ be the number of integers less than $n$ and prime to it, then if $x$ be any number less than $n$ and prime to it, we have $x^N=1$ (mod.\ $n$).   But we have in this case no analogue of a prime root---there is no number $x$, such that its several powers $x',\,x'',\,\ldots,\,x^{N-1}$ (mod.\ $n$) are all different from unity; or, what is the same thing, there is for each value of $x$ some submultiple of $N$, say $N'$, such that $x^{N'}=1$ (mod.\ $n$).   And these several numbers $N'$ have a least common multiple $I$, which is not $=N$, but is a submultiple of $N$; and this being so, then for all the several values of $x$, $I$ is said to be the maximum indicator.   For instance, $n=12$, $N=\totient(n)$; the numbers less than $12$ and prime to it are $1,\,5,\,7,\,11$.   We have, (mod.\ $12$), $1^1=1$, $5^2=1$, $7^2=1$, $11^2=1$, or the values of $N'$ are $1,\,2,\,2,\,2$; their least common multiple is $2$, and we have accordingly $I=2$: viz.\ $x^2=1$ (mod.\ $12$) has the $\totient(12)$ roots $1,\,5,\,7,\,11$.   So $n=24$, $\totient(n)=8$; the maximum indicator $I$ is in this case also $=2$.

\bigskip
\noindent\textbf{[p.~475]}\quad
A table of the maximum indicator $n=1$ to $1000$ is given by

Cauchy, \textit{Exer.\ d'Analyse} \&c., t.~ii.\ (1841), pp.~36--40, contained in the ``M\'emoire sur la r\'esolution des Equations ind\'etermin\'ees du premier degr\'e en nombres entiers,'' pp.~1--40.

\textbf{24.}\ \ It thus appears that for a composite number $n$, the $\totient(n)$ numbers less than $n$ and prime to it cannot be expressed as $\;\equiv\;$ (mod.\ $n$) to the power of a single root; but for the expression of them it is necessary to employ two or more roots.   A small table, $n=1$ to $50$, is given by

Cayley, ``Specimen Table $M\equiv a^\alpha b^\beta$ (mod.\ $N$) for any prime or composite modulus,'' \textit{Quart.\ Math.\ Journ.}\ vol.~ix.\ (1868), pp.~95, 96, and folding sheet, [397].

A specimen is
\[
\begin{array}{c|c}
\textit{Nos.}    & 12 \\
\textit{roots}   & 5,\,7 \\
\textit{Ind.}    & 2,\,2 \\
\textit{M.I.}    & 2 \\
\totient(\;)     & 4 \\\hline
 1 & 0,\,0 \\
 2 & \\
 3 & \\
 4 & \\
 5 & 1,\,0 \\
 6 & \\
 7 & 0,\,1 \\
 8 & \\
 9 & \\
10 & \\
11 & 1,\,1
\end{array}
\]
viz.\ for the modulus $12$ the roots are $5,\,7$, having respectively the indicators $2,\,2$, viz.\ $5^2=1$ (mod.\ $12$), $7^2=1$ (mod.\ $12$).   Hence also the maximum indicator is $=2$.   $\totient(=\totient(n))=4$ is the number of integers less than $12$ and prime to it, viz.\ these are $1,\,5,\,7,\,11$, which in terms of the roots $5,\,7$ and to mod.\ $12$ are respectively $=5^0.7^0,\ 5^1.7^0,\ 5^0.7^1$, and $5^1.7^1$.

\textbf{25.}\ \ \textit{Quadratic Residues.}\ ---In regard to a given prime number $p$, a number $N$ is or is not a quadratic residue according as the index of $N$ is even or odd, viz.\ $g$ being a prime root and $N=g^\alpha$, then $N$ is or is not a quadratic residue according as $\alpha$ is even or odd.   But the quadratic residues can, of course, be obtained directly without the consideration of prime roots.

A small table, $p=3$ to $97$ and $N=-1$ and (prime values) $3$ to $97$, is given by

Gauss, ``Disquisitiones Arithmeticae,'' 1801; Table II.\ (\textit{Werke}, t.~i.\ p.~469): I notice here a misprint in the top line of the original; it should be $-1,\,+2,\,+3,$ \&c., instead of

\bigskip
\noindent\textbf{[p.~476]}\quad
$-1,\,+2,\,+3,$ \&c.; the $-1$ is printed correctly on p.~499 of the French translation \textit{Recherches Arithm\'etiques}, Paris, 1807 and on p.~469 of vol.~I.\ of \textit{Werke}, (G\"ottingen, 1870).

A specimen is
\[
\begin{array}{c|cccccccccc}
   & -1 & +2 & +3 & +5 & +7 & +11 & +13 & +17 & +19 & +23\;\&c. \\\hline
19 &    & -  & -  &    &    &     & -   & -   & -   & -
\end{array}
\]
viz.\ $-1,\,2,\,3,\,13$ are not, $5,\,7,\,11,\,17$ \&c.\ are, quadratic residues of $19$.   The residues taken positively and less than $19$ are, in fact, $1,\,4,\,5,\,6,\,7,\,9,\,11,\,16,\,17$.

The same table carried from $p=3$ to $503$, and prime values $N=3$ to $997$, is given by

Gauss, \textit{Werke}, t.~ii.\ pp.~400--409.   A specimen is
\[
\begin{array}{c|cccccccccc}
   & 2 & 3 & 5 & 7 & 11 & 13 & 17 & 19 & 23\;\&c. \\\hline
19 & - &   & - &   &    & -  & -  & -  &
\end{array}
\]
viz.\ the arrangement is the same, except only that the $-1$ column is omitted.

\textbf{26.}\ \ We have also by Gauss

``Disquisitiones Arithmeticae'' Table III.\ (\textit{Werke}, t.~i.\ p.~470), for the conversion into decimals of a vulgar fraction, denominator $p$ or $p^\mu$, not exceeding $100$.   The explanation is given in Art.\ $314$ \textit{et seq.}\ of the same work.

But this table, carried to a greater extent, is given by Gauss, \textit{Werke}, t.~ii.\ pp.~412--434, ``Tafel zur Verwandlung gemeiner Br\"uche mit Nennern aus dem ersten Tausend in Decimalbr\"uche;''\ viz.\ the denominators are here primes or powers of primes, $p^\mu$ up to $997$.

To explain the table, consider a modulus $p^\mu$ (where $\mu$ may be $=1$); if $10$ is not a prime root of $p^\mu$, consider a prime root $r$, which is such that $r^e=10$ (mod.\ $p^\mu$), $e$ being a submultiple of $p^{\mu-1}(p-1)$; say we have $ef=p^{\mu-1}(p-1)$: then $10^f=1$ (mod.\ $p^\mu$).   Consider any fraction $\dfrac{N}{p^\mu}$; then we may write $N=r^{ke+l}$ (mod.\ $p^\mu$), $k$ from $0$ to $f-1$ and $l$ from $0$ to $e-1$, $N=10^k r^l$, and consequently $\dfrac{N}{p^\mu}$ and $\dfrac{10^k r^l}{p^\mu}$ have the same mantissa (decimal part regarded as an integer); hence, in order to know the mantissa of every fraction whatever of $\dfrac{N}{p^\mu}$, it is sufficient to know the mantissa of $\dfrac{r^l}{p^\mu}$, that is, the mantissas of $\dfrac{1}{p^\mu},\,\dfrac{r}{p^\mu},\,\dfrac{r^2}{p^\mu},\ldots,\dfrac{r^{e-1}}{p^\mu}$, or, what is the same thing, the mantissas of $\dfrac{10}{p^\mu},\,\dfrac{10r}{p^\mu},\,\ldots,\dfrac{10r^{e-1}}{p^\mu}$.

For instance, $p^\mu=11$, $10^2=1$ (mod.\ $11$), whence $f=2$, $e=5$; and taking $r=2$, we have $10=r^?$ (mod.\ $11$).

\bigskip
\noindent\textbf{[p.~477]}\quad
The required mantiss\ae, denoted in the table by
\[
(0),\;(1),\;(2),\;(3),\;(4),
\]
are those of
\[
\frac{10}{11},\;\frac{10.2}{11},\;\frac{10.2^2}{11},\;\frac{10.2^3}{11},\;\frac{10.2^4}{11},
\]
viz.\ these fractions are respectively
\[
\begin{aligned}
(0),\;(1),\;(2),\;(3),\;(4),\\
\cdot 9090\ldots,\ 1\cdot 8181\ldots,\ 3\cdot 6363\ldots,\ 7\cdot 2727\ldots,\ 14\cdot 5454\ldots,
\end{aligned}
\]
or their mantiss\ae\ are $90,\,81,\,63,\,27,\,54$.

And we accordingly have as a specimen
\[
11\ |\ (1)\ldots 81,\ (2)\ldots 63,\ (3)\ldots 27,\ (4)\ldots 54,\ (0)\ldots 90.
\]

Or again, as another specimen, $r=2$
\[
27\ |\ (1)\ldots 740,\ (2)\ldots 481,\ (3)\ldots 962,\ (4)\ldots 925,\ (5)\ldots 851,\ (0)\ldots 370.
\]

The table in this form extends to $p^\mu=463$; the values of $r$ (not given in the body of the table) are annexed, p.~420.

In the latter part of the table $p^\mu=467$ to $997$, we have only the mantiss\ae\ of $\dfrac{100}{p^\mu}$.   A specimen is
\[
\begin{array}{c|cccc}
547 & 1828153564 & 8994515539 & 3053016453 & 3820840950 \\
    & 6398537477 & 1480804387 & 5685557586 & 8372943327 \\
    & 2394881170 & 0182815356 & &
\end{array}
\]
viz.\ the fraction $\tfrac{100}{547}=\cdot 182815\ldots$ has a period of $91$, $=\tfrac{1}{6}.546$, figures.

\bigskip
\begin{center}
[F.~13.\ The Pellian Equation.]\ \ Art.~III.
\end{center}

\textbf{27.}\ \ The Pellian equation is $y^2=ax^2+1$, $a$ being a given integer number, which is not a square (or rather, if it be, the only solution is $y=1$, $x=0$), and $x,\,y$ being numbers to be determined: what is required is the least values of $x,\,y$, since these, being known, all other values can be found.   A small table $a=2$ to $68$ is given by

Euler, \textit{Op.\ Arith.\ Coll.}\ t.~i.\ p.~8.   The Memoir is ``Solutio problematum Diophanteorum per numeros integros,'' pp.~4--10, 1732--33.   The form of the table is
\[
\begin{array}{c|cc}
a & x\,(=p) & y\,(=q) \\\hline
2 & 2 & 3 \\
3 & 1 & 2 \\
5 & 4 & 9 \\
\vdots & \vdots & \vdots \\
68 & 4 & 33.
\end{array}
\]

\bigskip
\noindent\textbf{[p.~478]}\quad
Even here, for some of the values of $a$, the values of $x,\,y$ are extremely large; thus $a=61$, $x=226{,}153{,}980$, $y=1{,}766{,}319{,}049$.

And probably tables of a like extent may be found elsewhere; in particular, a table of the solution of $y^2=ax^2\pm 1$ ($-$ when the value of $a$ is such that there is a solution of $y^2=ax^2-1$, and $+$ for other values of $a$), $a=2$ to $135$, is given by Legendre, \textit{Th\'eorie des Nombres}, 2nd ed.\ 1808, in the Table X.\ (one page) at the end of the work.   For the before-mentioned number $61$, the equation is $y^2=61x^2-1$, and the values are $x=3805$, $y=29718$; much smaller than Euler's values for the equation $y^2=61x^2+1$.

\textbf{28.}\ \ The most extensive table, however, is given by

Degen, \textit{Canon Pellianus, sive Tabula simplicissimam \ae quationis celebratissim\ae\ $y^2=ax^2+1$ solutionem, pro singulis numeri dati valoribus ab $1$ ad $1000$ in usque numeris rationalibus, iisdemque integris exhibens.   Auctore Carolo Ferdinando Degen.   Hafni\ae\ (Copenhagen) apud Gerhardum Bonnarum, 1817.}\ \ 8vo.\ pp.\ iv to xxiv and $1$ to $112$.

The first table (pp.~3--106) is entitled as ``Tabula I.\ Solutionem Equationis $y^2-ax^2-1=0$ exhibens.''   It, in fact, also gives the expression of $\sqrt{a}$ as a continued fraction; thus a specimen is
\[
\begin{array}{c|ccccc}
209 & 14 & 2 & 5 & 3 & (2) \\
    & 1 & 13 & 5 & 8 & 11 \\
    & 3220 & & & & \\
    & 46551 & & & &
\end{array}
\]

Here the first line gives the continued fraction, viz.
\[
\sqrt{209}=14+\cfrac{1}{2+\cfrac{1}{5+\cfrac{1}{3+\cfrac{1}{2+\cfrac{1}{3+\cfrac{1}{5+\cfrac{1}{2+\cfrac{1}{28+\cfrac{1}{2+\&c.}}}}}}}}},
\]
the period being $(2,5,3,2,3,5,2)$ indicated by $2,5,3\,(2)$.   [The number of terms in the period is here odd, but it may be even; for instance, the period $(1,1,5,5,1,1)$ is indicated by $1,1\,(5,5)$.]

The second line contains auxiliary numbers presenting themselves in the process; thus, if $R^2=209$, we have $R=14+\dfrac{1}{\beta}$,
\[
\beta=\dfrac{1}{R-14} = \dfrac{1(R+14)}{209-14^2} = \dfrac{R+14}{13} = 2+\dfrac{1}{\beta'},
\]
\[
\beta'=\dfrac{13}{R-12} = \dfrac{13(R+12)}{209-12^2} = \dfrac{R+12}{5} = 5+\dfrac{1}{\gamma},
\]
\[
\gamma=\dfrac{5}{R-13} = \dfrac{5(R+13)}{209-13^2} = \dfrac{R+13}{8} = 3+\dfrac{1}{\delta},
\]
\&c.



\clearpage
\appendix
\section*{Original scan reference pages}
\noindent\footnotesize The following pages reproduce the corresponding range of the Internet Archive scan for visual reference of figures, tables, and any text that may have been compressed in the typeset reconstruction. Source: \texttt{collectedmathema09cayluoft.pdf}, pages 476--500.
\normalsize
\includepdf[pages=476-500,scale=0.92,pagecommand={\thispagestyle{empty}}]{local workspace path

\end{document}
